% 地支三合半合图
% 用途：展示四个三合局和对应的半合关系
\documentclass[tikz,border=10pt]{standalone}
\usepackage{ctex}
\usepackage{xcolor}
\pagecolor{white}
\usetikzlibrary{arrows.meta, positioning, calc, decorations.pathmorphing, backgrounds}

\begin{document}
\begin{tikzpicture}[scale=1.2]
  \coordinate (zi)  at ( 90:4.0);
  \coordinate (chou) at ( 60:4.0);
  \coordinate (yin)  at ( 30:4.0);
  \coordinate (mao)  at (  0:4.0);
  \coordinate (chen) at (330:4.0);
  \coordinate (si)   at (300:4.0);
  \coordinate (wu)   at (270:4.0);
  \coordinate (wei)  at (240:4.0);
  \coordinate (shen) at (210:4.0);
  \coordinate (you)  at (180:4.0);
  \coordinate (xu)   at (150:4.0);
  \coordinate (hai)  at (120:4.0);

  \draw[thick, gray!60] (0,0) circle (4.0);

  \foreach \name/\angle/\text in {
    zi/90/子, chou/60/丑, yin/30/寅, mao/0/卯,
    chen/330/辰, si/300/巳, wu/270/午, wei/240/未,
    shen/210/申, you/180/酉, xu/150/戌, hai/120/亥
  } {
    \fill[black] (\angle:4.0) circle (2pt);
    \node[font=\Large\bfseries] at (\angle:4.8) {\text};
  }

  % 核心地支（子卯午酉）高亮
  \draw[orange!80!black, ultra thick] (90:4.0)  circle (12pt);
  \draw[teal!80!black,  ultra thick] (0:4.0)   circle (12pt);
  \draw[red!80!black,   ultra thick] (270:4.0) circle (12pt);
  \draw[gray!80!black,  ultra thick] (180:4.0) circle (12pt);
  \node[font=\Large\bfseries, orange!80!black] at (90:4.0)  {子};
  \node[font=\Large\bfseries, teal!80!black]  at (0:4.0)   {卯};
  \node[font=\Large\bfseries, red!80!black]   at (270:4.0) {午};
  \node[font=\Large\bfseries, gray!80!black]  at (180:4.0) {酉};

  % 三合局：实线三角 + 标签
  \draw[orange!70!black, very thick] (shen) -- (zi) -- (chen) -- cycle;
  \node[font=\normalsize\bfseries, orange!70!black, fill=white, inner sep=2pt] at (120:1.8) {水局};
  \draw[teal!70!black, very thick] (hai) -- (mao) -- (wei) -- cycle;
  \node[font=\normalsize\bfseries, teal!70!black, fill=white, inner sep=2pt] at (60:2.0) {木局};
  \draw[red!70!black, very thick] (yin) -- (wu) -- (xu) -- cycle;
  \node[font=\normalsize\bfseries, red!70!black, fill=white, inner sep=2pt] at (300:2.0) {火局};
  \draw[gray!60!black, very thick] (si) -- (you) -- (chou) -- cycle;
  \node[font=\normalsize\bfseries, gray!60!black, fill=white, inner sep=2pt] at (240:1.8) {金局};

  % 半合：虚线连接核心地支与两边
  \draw[orange!70!black, dashed, thin] (zi) -- (shen);
  \draw[orange!70!black, dashed, thin] (zi) -- (chen);
  \draw[teal!70!black, dashed, thin] (mao) -- (hai);
  \draw[teal!70!black, dashed, thin] (mao) -- (wei);
  \draw[red!70!black, dashed, thin] (wu) -- (yin);
  \draw[red!70!black, dashed, thin] (wu) -- (xu);
  \draw[gray!60!black, dashed, thin] (you) -- (si);
  \draw[gray!60!black, dashed, thin] (you) -- (chou);

  % 图例
  \node[font=\small\bfseries, anchor=west] at (-6.0,-5.5) {\textbf{三合局}（实线三角）};
  \node[font=\small, anchor=west] at (-6.0,-6.0) {橙:水   青:木   红:火   灰:金};
  \node[font=\small\bfseries, anchor=west] at (-6.0,-6.5) {\textbf{半合}（虚线）：聚合倾向，不等于完整成局};

  % 白色背景，防止 GitHub dark mode 下透明 SVG 看不清
  \begin{scope}[on background layer]
    \fill[white] (current bounding box.south west) rectangle (current bounding box.north east);
  \end{scope}
\end{tikzpicture}
\end{document}
